\chapter{Domain Protocol Cards and Transfer Obligations}
\label{app:domain-protocol-cards}

Each domain chapter follows one template so that universality is tested through comparable protocol cards rather than loose analogy. This appendix collects those cards in one place for defense preparation and future replication.

\begin{table}[t]
\centering
\begin{tabular}{p{0.24\textwidth}p{0.66\textwidth}}
\toprule
\textbf{Theorem family} & \textbf{Instantiation question answered by the appendix}\\
\midrule
T1 & What is the true-state violation event in the domain?\\
T2 & How is the observation surface separated from the true physical state?\\
T3 & What admissible action geometry and repair operator are used?\\
T4 & How is the certificate payload mapped into domain semantics?\\
T5--T8 & What temporal, behavioral, and transfer obligations remain open?\\
\bottomrule
\end{tabular}
\caption{Domain instantiation summary across the ORIUS theorem families.}
\label{tab:monograph-domain-instantiation-summary}
\end{table}

\section{Battery Energy Storage}
\label{app:protocol-battery}
\begin{longtable}{p{0.24\textwidth}p{0.68\textwidth}}
\caption{Protocol card for Battery Energy Storage.}\\
\toprule
\textbf{Field} & \textbf{Definition}\\
\midrule
Evidence tier & reference\\
System context & Grid-scale storage dispatch under delayed, stale, or dropped telemetry can appear feasible on observed state while violating true state-of-charge limits in the physical asset.\\
Safety predicate & The defended predicate is zero true-state violation of battery SOC and dispatch feasibility under degraded observation, with bounded intervention cost rather than unconstrained throughput maximization.\\
Adapter mapping & The battery adapter binds telemetry parsing, reliability scoring, uncertainty inflation, action tightening, and certificate emission to energy dispatch while preserving the universal Detect-Calibrate-Constrain-Shield-Certify contract.\\
Telemetry degradation model & Faults include dropout, staleness, delayed measurements, and spoof-like anomalies that widen the observation-action safety gap by masking the true plant state.\\
Dataset and protocol & The reference evidence is built from tracked German and US energy artifacts, battery replay traces, and the locked publication tables that already support the canonical ORIUS claim family.\\
Results & In the locked reference surface, ORIUS closes the measured true-state violation rate from 3.90 percent to 0.00 percent while keeping the intervention rate materially below a shutdown-only fallback policy.\\
Limitations & Battery is still the only row with full theorem-grade depth. It should therefore anchor the strongest proof language without being allowed to dominate the conceptual framing of the full monograph.\\
Transfer obligations & Other domains must match battery on telemetry realism, replay discipline, certificate semantics, and closure under the universal benchmark before they can claim equal validation depth.\\
Promotion gate & retain theorem-grade witness status while serving as the calibration surface for new domains\\
\bottomrule
\end{longtable}

The transfer obligation for this domain is not merely to preserve code portability. It is to preserve the semantic contract between degraded observation, action repair, and certificate emission under the same benchmark discipline used everywhere else in the book.

\section{Autonomous Vehicles}
\label{app:protocol-av}
\begin{longtable}{p{0.24\textwidth}p{0.68\textwidth}}
\caption{Protocol card for Autonomous Vehicles.}\\
\toprule
\textbf{Field} & \textbf{Definition}\\
\midrule
Evidence tier & proof\_validated\\
System context & Longitudinal autonomy must preserve headway, speed, and collision margins while acting through delayed or degraded perception channels that can misstate closing speed or lead-vehicle behavior.\\
Safety predicate & The defended predicate is the absence of true-state spacing or speed violations after the repaired control action is executed on the vehicle model.\\
Adapter mapping & The AV adapter maps tracked kinematic telemetry into the common ORIUS contract and translates the repaired action into a constrained acceleration command.\\
Telemetry degradation model & The dominant faults are stale leader observations, burst dropout, unrealistic spikes, and delayed kinematic packets that distort the observed stopping margin.\\
Dataset and protocol & The current row is backed by locked trajectory telemetry and a shared universal fault schedule under the same replay harness used by the other non-battery domains.\\
Results & Under the current TTC-based safety surface and predictive-entry-barrier formulation, the locked replay closes the residual true-state violation surface materially enough to clear the bounded proof-validation gate for the present longitudinal setting.\\
Limitations & The row is still bounded to the current longitudinal autonomy contract. Multi-lane interaction, richer state repair, and broader deployment semantics remain outside the present validated surface.\\
Transfer obligations & Future promotion work must preserve the single-contract TTC semantics, replay-backed soundness checks, and explicit fallback compatibility instead of reopening multiple competing repair formulations.\\
Promotion gate & retain proof-validated status while extending TTC closure to richer vehicle interaction settings\\
\bottomrule
\end{longtable}

The transfer obligation for this domain is not merely to preserve code portability. It is to preserve the semantic contract between degraded observation, action repair, and certificate emission under the same benchmark discipline used everywhere else in the book.

\section{Industrial Process Control}
\label{app:protocol-industrial}
\begin{longtable}{p{0.24\textwidth}p{0.68\textwidth}}
\caption{Protocol card for Industrial Process Control.}\\
\toprule
\textbf{Field} & \textbf{Definition}\\
\midrule
Evidence tier & proof\_validated\\
System context & Industrial plants face thermal, pressure, and power safety constraints under unreliable sensor networks where stale or missing observations can hide excursions until they have already become operational hazards.\\
Safety predicate & The defended predicate is zero true-state violation of the bounded process envelope after the repaired action is issued, while preserving plant continuity better than a brute shutdown response.\\
Adapter mapping & The industrial adapter reuses the universal kernel by swapping in process-specific state features, constraint evaluation, and action repair against the plant envelope.\\
Telemetry degradation model & The row stresses packet loss, stale reads, spikes, and out-of-order process measurements that mimic degraded plant telemetry rather than laboratory noise.\\
Dataset and protocol & The current surface uses the locked industrial telemetry row, the verified training audit, and the shared universal replay harness.\\
Results & On the tracked evidence row, ORIUS reduces true-state violations from 21.53 percent to 0.00 percent, making industrial control one of the strongest non-battery demonstrations of universality in practice.\\
Limitations & The proof-validated label is still bounded: the row is a defended domain instantiation, not a claim that every industrial control topology inherits the same guarantee automatically.\\
Transfer obligations & Future industrial extensions must retain real telemetry, replay closure, and explicit certificate semantics rather than relying on architecture-only analogies.\\
Promotion gate & retain proof-validated status with broader plant families and stronger controller diversity\\
\bottomrule
\end{longtable}

The transfer obligation for this domain is not merely to preserve code portability. It is to preserve the semantic contract between degraded observation, action repair, and certificate emission under the same benchmark discipline used everywhere else in the book.

\section{Medical and Healthcare Monitoring}
\label{app:protocol-healthcare}
\begin{longtable}{p{0.24\textwidth}p{0.68\textwidth}}
\caption{Protocol card for Medical and Healthcare Monitoring.}\\
\toprule
\textbf{Field} & \textbf{Definition}\\
\midrule
Evidence tier & proof\_validated\\
System context & Medical telemetry can fail through alarm dropout, stale vitals, or delayed charting, allowing observed safety to diverge from the patient state that should govern escalation or suppression decisions.\\
Safety predicate & The defended predicate is zero true-state violation of the monitored physiologic safety envelope under the bounded intervention semantics carried by the current healthcare adapter.\\
Adapter mapping & The healthcare adapter binds physiologic telemetry, uncertainty widening, threshold-aware action repair, and certificate logging to the common ORIUS kernel without changing the underlying runtime contract.\\
Telemetry degradation model & Faults include stale vital-sign packets, missing waveform updates, spikes, and delayed observations that can suppress or delay clinically necessary intervention.\\
Dataset and protocol & The row uses the locked ICU vitals evidence surface and the shared universal replay protocol.\\
Results & The present locked row reduces true-state violations from 6.25 percent to 0.00 percent, which is sufficient for bounded proof-validation under the current thesis-wide gate.\\
Limitations & This is not a claim of universal clinical deployment. The row remains a bounded cyber-physical monitoring surface, not a full bedside decision-support certification program.\\
Transfer obligations & Future medical promotion requires stronger outcome semantics, wider patient cohorts, and explicit regulatory traceability beyond the current replay-centered scope.\\
Promotion gate & multi-cohort clinical validation and richer intervention outcome accounting\\
\bottomrule
\end{longtable}

The transfer obligation for this domain is not merely to preserve code portability. It is to preserve the semantic contract between degraded observation, action repair, and certificate emission under the same benchmark discipline used everywhere else in the book.

\section{Navigation and Guidance}
\label{app:protocol-navigation}
\begin{longtable}{p{0.24\textwidth}p{0.68\textwidth}}
\caption{Protocol card for Navigation and Guidance.}\\
\toprule
\textbf{Field} & \textbf{Definition}\\
\midrule
Evidence tier & shadow\_synthetic\\
System context & Navigation systems must preserve corridor, obstacle, or path-feasibility constraints even when localization and guidance telemetry degrade under dropout or stale updates.\\
Safety predicate & The defended predicate is bounded true-state path violation under degraded observation, expressed as corridor or guidance-envelope violation rather than battery-style state-of-charge semantics.\\
Adapter mapping & The navigation adapter maps localization and path telemetry to the universal kernel, but the current row remains a portability and closure surface rather than a real-data proof surface.\\
Telemetry degradation model & The current row includes synthetic dropout, stale localization, corrupted headings, and delayed guidance packets.\\
Dataset and protocol & The row is intentionally flagged as synthetic in the closure matrix and is treated as portability evidence rather than parity evidence.\\
Results & ORIUS reduces the synthetic-row TSVR materially, but not to zero, which is why the book treats navigation as an architecture-validating row rather than a defended empirical peer.\\
Limitations & No claim of field-grade navigation validation should be made from the current row.\\
Transfer obligations & Promotion requires locked real-data navigation telemetry, verified replay surfaces, and the same certificate semantics enforced elsewhere in the book.\\
Promotion gate & locked real-data navigation telemetry and replay-backed closure\\
\bottomrule
\end{longtable}

The transfer obligation for this domain is not merely to preserve code portability. It is to preserve the semantic contract between degraded observation, action repair, and certificate emission under the same benchmark discipline used everywhere else in the book.

\section{Aerospace Control}
\label{app:protocol-aerospace}
\begin{longtable}{p{0.24\textwidth}p{0.68\textwidth}}
\caption{Protocol card for Aerospace Control.}\\
\toprule
\textbf{Field} & \textbf{Definition}\\
\midrule
Evidence tier & experimental\\
System context & Aerospace safety layers must preserve envelope constraints under degraded flight-state observation, delayed telemetry, and actuator-response uncertainty.\\
Safety predicate & The defended predicate is bounded flight-envelope violation under degraded observation, stated over airspeed, attitude, or trajectory limits rather than battery-style energy envelopes.\\
Adapter mapping & The aerospace adapter proves that the ORIUS contract is expressive enough for flight-state envelopes, but the row is still a systems-experimental surface rather than a mature empirical benchmark.\\
Telemetry degradation model & The current row models delayed and degraded flight-state telemetry with bounded envelope checks.\\
Dataset and protocol & The row uses locked replay artifacts, but the present evidence is insufficient for promotion because the repaired action does not yet close the violation gap materially.\\
Results & The tracked status table keeps aerospace at the experimental tier because the current row does not improve the governing TSVR metric enough to satisfy the universal promotion gate.\\
Limitations & The row should be used to discuss scope, not to overstate generality.\\
Transfer obligations & Promotion requires stronger plant realism, better envelope repair, and a nontrivial post-repair gain under the tracked protocol.\\
Promotion gate & material post-repair improvement on a stronger flight-envelope benchmark\\
\bottomrule
\end{longtable}

The transfer obligation for this domain is not merely to preserve code portability. It is to preserve the semantic contract between degraded observation, action repair, and certificate emission under the same benchmark discipline used everywhere else in the book.

